% Push streaming on D2Q9 and the periodic wrap of torch.roll.
\begin{tikzpicture}[
  >={Stealth[length=2.2mm]},
  node/.style={circle, fill=gray!25, draw=gray!60, inner sep=0pt, minimum size=6pt},
  pop/.style={->, line width=0.9pt, blue!60!black},
  hi/.style={->, line width=1.5pt, red!70!black},
  lbl/.style={font=\small, inner sep=1pt},
  note/.style={font=\small, gray!80!black, align=left},
  cell/.style={draw, fill=blue!7, minimum size=9mm, font=\small\ttfamily},
]
% ---- left: push from one node to its eight neighbours ----
\begin{scope}[shift={(0,0)}]
  \draw[gray!30, step=2.4] (-2.5,-2.5) grid (2.5,2.5);
  \foreach \x in {-2.4,0,2.4} \foreach \y in {-2.4,0,2.4} \node[node] at (\x,\y) {};
  \node[node, fill=black!70] (x) at (0,0) {};
  \node[lbl, below left=2pt] at (x) {$\mathbf{x}$};
  % seven ordinary populations leave the node; label sits beside the arrow tip
  \foreach \i/\dx/\dy in {2/0/1, 3/-1/0, 4/0/-1, 5/1/1, 6/-1/1, 7/-1/-1, 8/1/-1} {
    \draw[pop] (0,0) -- ({2.4*\dx - 0.3*\dx}, {2.4*\dy - 0.3*\dy});
    \node[lbl, blue!60!black] at ({1.55*\dx + 0.42*\dy}, {1.55*\dy - 0.42*\dx}) {$f_{\i}^{*}$};
  }
  % the highlighted one: f_1^*(x) becomes f_1(x + c_1)
  \draw[hi] (0,0) -- (2.1,0);
  \node[lbl, red!70!black, above] at (1.1,0.08) {$f_1^{*}(\mathbf{x},t)$};
  \node[node, fill=red!70!black] at (2.4,0) {};
  \node[lbl, red!70!black, align=center, anchor=north] at (2.4,-0.25)
    {$f_1(\mathbf{x}+\vc_1\Delta t,\, t+\Delta t)$};
  % rest population stays
  \node[lbl, black!70, anchor=north east] at (-0.15,-0.35) {$f_0^{*}$ stays};
  \node[font=\small\bfseries, align=center] at (0,-4.0) {push: every $f_i^{*}$ leaves along its own $\vc_i$};
\end{scope}

% ---- right: one direction on one axis, periodic wrap ----
\begin{scope}[shift={(6.3,1.0)}]
  \node[note, anchor=south west, font=\small\ttfamily] at (-0.4,1.05) {torch.roll(f[1], shifts=(+1, 0), dims=(0, 1))};
  % before
  \foreach \k in {0,...,4} {
    \node[cell] (a\k) at (\k*1.0,0.4) {\k};
  }
  \node[lbl, left=4pt of a0] {$f_1^{*}$};
  \node[lbl, right=4pt of a4] {$x \rightarrow$};
  % after
  \foreach \k in {0,...,4} {
    \pgfmathtruncatemacro{\src}{mod(\k+4,5)}
    \node[cell] (b\k) at (\k*1.0,-1.6) {\src};
  }
  \node[lbl, left=4pt of b0] {$f_1$};
  \foreach \k in {0,...,3} {
    \pgfmathtruncatemacro{\kk}{\k+1}
    \draw[pop] (a\k.south) -- (b\kk.north);
  }
  % wrap: last cell goes to first
  \draw[hi, rounded corners=3pt] (a4.south) -- ++(0,-0.45) -- ++(-4.0,0) -- (b0.north);
  \node[lbl, red!70!black, anchor=west] at (4.55,-0.55) {wrap};
  \node[note, anchor=north west, text width=6.4cm] at (-0.4,-2.35)
    {Cell $x$ holds the value that came from $x-1$; the value leaving the last cell re-enters
     at the first. A wall boundary overwrites that entry afterwards.};
  \node[font=\small\bfseries] at (2.0,-5.0) {periodic by construction};
\end{scope}
\end{tikzpicture}
